The source assumptions ensure that the supplied representation and every monomially rescaled disturbance system remain within the stated source-minimal non-Gaussian fixed-law setting; the algebraic sharpness argument itself uses the full-row-rank and solvability conditions explicitly.

First, \(J_x(C)\) is nonempty. Extend \(C\) to \(T_0\) and write \(W_0=T_0^{-1}\), as in the witness construction. The \((x,x)\) entry of \(T_0W_0=I_n\) is
\[
 1=\sum_{j=1}^n C_{xj}(W_0)_{jx}. \tag{10}
\]
Thus some \(j\) satisfies \(C_{xj}(W_0)_{jx}\ne0\). The inverse-cofactor formula then implies \(\det((T_0)_{-x,-j})\ne0\). The observed rows other than \(x\), restricted away from column \(j\), are a subset of the rows of this nonsingular minor, so \(\operatorname{rank}(C_{-x,-j})=p-1\). Together with \(C_{xj}\ne0\), this proves \(j\in J_x(C)\).

For the first set inclusion, take any \(\mathcal M\in\mathfrak F_x(C,\varepsilon)\). The necessary-deletion-rank lemma gives \(s_{\mathcal M}(x)\in J_x(C)\), and the intervention-ratio lemma gives
\[
 \theta_{y\leftarrow x}(\mathcal M)
 =\frac{C_{y,s_{\mathcal M}(x)}}{C_{x,s_{\mathcal M}(x)}}
 \in\mathcal R_{xy}(C). \tag{11}
\]
Conversely, if \(r_0\in\mathcal R_{xy}(C)\), then by definition \(r_0=C_{yj}/C_{xj}\) for some \(j\in J_x(C)\). The shear-and-matching sufficiency lemma returns \(\mathcal M_j\in\mathfrak F_x(C,\varepsilon)\) with assigned source \(j\), and the intervention-ratio lemma yields \(\theta_{y\leftarrow x}(\mathcal M_j)=r_0\). This proves the reverse inclusion and constructive sharpness.

For the arithmetic bound, one Gaussian elimination extends the \(p\) rows of \(C\) and inverts \(T_0\) in \(O(n^3)\) operations. Criterion (5) in the sufficiency proof classifies all \(j\) by scanning row \(j\) of \(W_0\), so all deletion decisions cost \(O(n^2)\) after the inverse. For each of at most \(n\) admissible indices, the shear costs \(O(n^2)\), a perfect matching in the \((n-1)\)-by-\((n-1)\) nonzero pattern can be found by an augmenting-path algorithm in \(O(n^3)\), and row normalization costs \(O(n^2)\). Hence all witnesses cost \(O(n^4)\) arithmetic operations. Exact zero tests are part of the real-arithmetic model; this is not a floating-point or bit-complexity claim. Selecting one representative \(j\) for each distinct ratio returns one exact witness per distinct value.